\topmargin=-1cm
\chapter*{Sommario}
Il crowdsensing mobile consente la raccolta distribuita di dati
geografici aggregando i contributi di un elevato numero di
partecipanti.
Questo paradigma si presta a scenari eterogenei come ad esempio: monitoraggio
ambientale, analisi del traffico urbano, raccolta di dati
epidemiologici.
Scenari in cui la granularità spaziale e la scalabilità
del numero di contributori sono requisiti primari.


Le soluzioni esistenti affidano il coordinamento della campagna a
un server centralizzato, introducendo un unico punto di fiducia
che può compromettere la riservatezza dei dati raccolti e la
privacy dei partecipanti.
Un server centralizzato può essere compromesso, può colludere con
terze parti o essere soggetto a guasti, rendendo fragile l'intero
sistema.

Questa tesi presenta il design, l'implementazione e la valutazione
di un protocollo di crowdsensing decentralizzato e rispettoso della
privacy, in cui il coordinamento è delegato a uno smart contract
Ethereum.


Lo smart contract sostituisce il server fidato con un intermediario
pubblicamente verificabile e resistente alle manomissioni: la sua
logica è eseguita deterministicamente dall'Ethereum Virtual
Machine (EVM), ogni nodo della rete ne mantiene una copia, e il
comportamento del contratto non può essere alterato da alcuna
singola parte dopo il deploy.

Il protocollo coinvolge tre attori principali.
Il campaign initiator, \texttt{B} avvia la campagna, definisce i parametri
spaziali e distribuisce le chiavi crittografiche.
I sender raccolgono e inviano i dati geografici.
I verifier, selezionati casualmente dal contratto, validano le
submission esprimendo un voto di approvazione o rifiuto.
La partecipazione di sender e verifier è regolata da un meccanismo
di cauzione tramite token ERC-20 personalizzato, \texttt{DST}: il deposito
viene restituito al termine della campagna, unitamente al reward,
solo ai partecipanti onesti.

La riservatezza dei dati è garantita attraverso crittografia
ibrida.
Ogni sender cifra il proprio payload con una chiave simmetrica
AES-256-GCM e cifra quest'ultima con la chiave pubblica RSA del
verifier eletto tramite RSA-OAEP.
Il testo cifrato viene archiviato off-chain su IPFS; on-chain viene
pubblicato esclusivamente il Content Identifier~(CID) cifrato.

L'area della campagna è suddivisa in una griglia spaziale di chunk:
il contratto valida on-chain l'origine geografica di ciascuna
submission e limita il numero di dati accettati per chunk,
prevenendo la sovra-rappresentazione delle zone ad alta densità di
partecipanti.

Un'analisi di sicurezza e incentivi del protocollo base individua
tre problemi aperti.
Primo, un verifier malevolo può accedere ai dati in chiaro prima di
esprimere il proprio verdetto, falsando la valutazione.
Secondo, la distribuzione geografica dei reward non incentiva i
sender a contribuire alle aree meno coperte, favorendo la
concentrazione nelle zone già sature.
Terzo, i verifier che esprimono voti errati non subiscono alcuna
conseguenza economica, oltre alla perdita del deposito cauzionale,
rendendo conveniente il comportamento scorretto.
La correttezza del protocollo base rispetto agli obiettivi di
segretezza e autenticazione è verificata formalmente con ProVerif
sotto il modello di attaccante Dolev-Yao, solo nel caso in cui gli attori si comportino onestamente e nel caso in cui l'avversario $\mathcal{A}$ non si autentichi come verifier.

Un secondo contratto Solidity, \texttt{PSC}, estende il protocollo base
affrontando ciascuno dei tre problemi.
Per il problema del verifier malevolo, PSC introduce un ulteriore
strato crittografico: \texttt{B}, per ogni dato inviato crea una coppia di chiavi asimmetriche.
Quindi \texttt{B} decifra il dato inviato, riesegue l'upload del dato cifrato su IPFS e successivamente attraverso la chiave pubblica dei verificatori eletti distribuisce la chiave di decrittazione del dato da verificare.

Per lo squilibrio geografico dei reward, PSC implementa una funzione
di decadimento basata sul rapporto di saturazione del chunk: il
reward per una submission verso un chunk già parzialmente saturo
viene ridotto secondo la formula~$(1 - \alpha)^{1/2^k}$, dove
$\alpha$~è il rapporto di saturazione e~$k$~è un parametro
configurabile che controlla la ripidità del decadimento.

La radice frazionaria viene calcolata on-chain tramite il metodo
babilonese iterativo, evitando operazioni in virgola mobile non
supportate dall'EVM.

Per l'assenza di penalità, PSC traccia i voti corretti e scorretti
di ciascun verifier e consente al campaign initiator di espellere
i verifier che superano una soglia configurabile di voti errati.

L'implementazione comprende i due contratti Solidity documentati,
un SDK TypeScript che espone un'API per-attore e gestisce le
operazioni crittografiche off-chain, e una suite di test completa
sviluppata con Hardhat.

Le ottimizzazioni di gas documentate includono la gestione della
pending list tramite swap-and-pop al posto dello shift lineare FIFO,
la compressione delle struct in un singolo slot EVM da 32~byte,
l'aritmetica \texttt{unchecked} nelle funzioni di reward, e
l'eliminazione dei dati di enrollment tramite \texttt{delete} dopo
il loro utilizzo.

La valutazione empirica copre tre dimensioni.
L'analisi dei costi di gas su 16~scenari, con sender e verifier
nell'intervallo~$\{1, 10, 100, 1000\}$, mostra che il contratto
ottimizzato riduce i costi del verifier del~98\% rispetto alla
baseline FIFO.
La baseline FIFO risulta non praticabile su reti Fusaka-era in tutti
i 16~scenari, superando il limite per-transazione di~$2^{24}$~gas
introdotto dall'EIP-7825 di un fattore~$3.4\times$.
PSC aggiunge un overhead contenuto rispetto a PFC, con l'incremento
più significativo sulla data submission dovuto al calcolo del reward decaduto.

Su reti L2 quali Arbitrum e Optimism, il costo di partecipazione
diventa trascurabile per tutti i ruoli.
L'analisi della saturazione dei blocchi sotto campagne concorrenti
evidenzia una sezione critica sulla capacità del blocco: con quattro
campagne parallele, il numero medio di blocchi consumati per
campagna raggiunge~$3.65\times$.
La simulazione spaziale sul dataset di traiettorie taxi della città
di Porto valida il meccanismo di decadimento del reward, mostrando
come la funzione redistribuisca efficacemente gli incentivi verso le
aree meno coperte.
La copertura geografica satura attorno al~50\% tra 500 e
1\,000~sender, un limite attribuibile alla concentrazione geografica
delle traiettorie nell'area urbana centrale piuttosto che al
protocollo.

I risultati confermano la fattibilità del protocollo su reti L2
e su blockchain a basso costo del gas.
Le direzioni future includono l'eliminazione del trust residuo nel
campaign initiator tramite Proxy Re-Encryption decentralizzata,
l'estensione della verifica formale a scenari di enrollment
disonesto, e la validazione su dataset con distribuzione geografica
più uniforme.